%% ======================================================================
%%  front matter: title page, copyright, contents, translator's
%%  introduction, introduction, notations
%% ======================================================================

%% ---------------------------- title page ------------------------------
\thispagestyle{empty}%
\begin{center}
\vspace*{14mm}
{\scshape Pierre Samuel}\\[2pt]
{\itshape\small Professor at the Faculty of Sciences, Paris}
\vspace{44mm}

{\fontsize{24}{29}\selectfont Algebraic Theory}\\[3mm]
{\fontsize{24}{29}\selectfont of}\\[3mm]
{\fontsize{24}{29}\selectfont Numbers}
\vspace{34mm}

{\bfseries\small Translated from the French by Allan J. Silberger}
\vfill
{\small HERMANN}\\[2pt]
{\small Publishers in Arts and Science, Paris, France}
\end{center}
\clearpage

%% -------------------------- copyright page ----------------------------
\thispagestyle{empty}%
\begingroup\small
\noindent
Translated from the original French text \emph{Théorie algébrique des
nombres}, published by Hermann, Paris, in 1967 in their series
\emph{Méthodes}.
\vfill
\noindent
Printed in Great Britain by William Clowes and Sons, Limited, London and
Beccles\\[4pt]
\copyright{} HERMANN, PARIS, 1970\\[4pt]
All rights reserved. No part of this work may be reproduced or transmitted
in any form or by any means, electronic or mechanical, including
photocopying and recording, or by any information storage or retrieval
system, without permission in writing from the original publisher.
\endgroup
\clearpage

%% ------------------------------ contents ------------------------------
\setcounter{page}{5}%
\renewcommand{\chaptertitle}{CONTENTS}%
\tableofcontents

%% ---------------------- translator's introduction ---------------------
\clearpage
\setcounter{page}{7}%
\fmhead{Translator's introduction}{TRANSLATOR'S INTRODUCTION}{translator}
\begin{flushright}\small\itshape
LE JUGE\quad Accusé, soyez fidèle à la vérité\\
L'ACCUSÉ\quad Mais je suis coupable d'avoir commis le crime
\end{flushright}
\vspace{1.5\baselineskip}

Professor Samuel has written a beautiful book and, in doing so, he has made
an outstanding contribution to the teaching of mathematics. The reader of
this book needs to know little beyond the content of a first course in
abstract algebra. He will find here a very clear account of classical
algebraic number theory, including sufficient examples, applications, and
exercises. This book will whet the reader's appetite for and provide him
with the background he needs to enjoy the lengthier treatises listed in the
bibliography.

The translator wishes to acknowledge, with gratitude, support and assistance
he has received while translating this book from the following sources:
Bowdoin College, the National Science Foundation (grants GP-7139 and
GP-9388), and especially Mrs Jean Hughes who typed the manuscript.

He also wishes to thank various readers of the earlier French edition,
especially Mrs. Germaine Revuz and Mr. Alain Bouvier, who have communicated
lists of errata. These mistakes have been corrected in the present English
edition.

\vspace{1.5\baselineskip}
\begin{flushright}
\textsc{Allan J. Silberger}\\
Brunswick, Maine, May 1, 1969
\end{flushright}

%% ---------------------------- introduction ----------------------------
\clearpage
\setcounter{page}{9}%
\fmhead{Introduction}{INTRODUCTION}{introduction}

The theory of numbers, or arithmetic, is often called ``the queen of
mathematics''. The simplicity of its subject matter (the ordinary integers
and their generalizations), the elegance and diversity of its methodology,
its numerous unsolved problems strongly attract mathematicians of all
classes, whether they be beginners, professional number theorists, or
specialists in other branches of mathematics. It should not greatly surprise
the reader of this book to learn that its author is an algebraic geometer,
who has published no research in what may properly be called arithmetic.

Thus, the reader should not expect an especially profound or comprehensive
treatment of number theory. Moreover, he will meet only one of many
points-of-view from which one may approach number theory, namely the
algebraic point-of-view. Aside from an elementary result of Minkowski
concerning lattices in $\mathbf{R}^{n}$, the beautiful and fertile technique
of analytic number theory is entirely omitted from this book.

Emphasis upon the algebraic point-of-view seems to me justifiable for
several reasons. First of all, the algebraic point-of-view establishes the
context in which number theoretic problems have their most natural
formulation. This is true of even those problems which concern only the
natural numbers. For example, the problem of finding all integer solutions
of the Pell-Fermat equation $X^{2}-dY^{2}=\pm1$ ($d$: a square-free integer)
involves in an essential way the study of the quadratic field
$\mathbf{Q}[\sqrt{d}]$. For the ``great'' equation of Fermat
$X^{n}+Y^{n}=Z^{n}$ the field of $n$th roots of unity plays an analogous
role. In order to represent an integer as the sum of two (respectively,
four) squares, it is advantageous to work in the ring of Gaussian integers
(respectively, in a suitably chosen quaternion algebra). The law of
quadratic reciprocity involves both quadratic fields and roots of unity.
Fields more general than the rational numbers and rings more general than
the ordinary integers arise quite naturally when one discusses any of the
above problems.

Secondly, although the algebraic approach does not lead to a solution of all
number theoretic problems, it does, as the reader will see, nonetheless
quickly lead to substantial results. Continuing in the direction of this
book, one would reach the deep theorems of class field theory.

Thirdly, even those who prefer analytic number theory will agree that the
full generality and power of the analytic approach reveals itself only in
the context of number fields and simple algebras, not in investigations
involving the rational numbers alone. For example, the Dedekind zeta
function of a number field possesses properties quite analogous to those of
the ordinary Riemann zeta function and it is not much more difficult to
establish these properties in complete generality. A similar statement holds
for the various $L$-series.

Finally, algebraic number theory provides the student with numerous
illustrative examples of notions he has encountered in his algebra courses:
groups, rings, fields, ideals, quotient rings and quotient fields,
homomorphisms and isomorphisms, modules and vector spaces. A further benefit
to the student lies in the fact that, in studying algebraic number theory,
he will meet many new algebraic notions, notions which are fundamental not
only for arithmetic but for other branches of mathematics as well, in
particular algebraic geometry. Here are some examples: integrality, field
extensions, Galois theory, modules over principal ideal rings, Noetherian
rings and modules, Dedekind rings, and rings of fractions.

The preceding describes implicitly what the reader will find in this book
and what he will not. I have assumed that the reader knows basic
algebra---elementary facts concerning groups, rings, polynomial rings, and
vector spaces---the basic facts concerning subobjects, quotient objects, and
product objects in these categories---how to pass to the quotient by an
ideal or submodule---the various different notions of homomorphism and
isomorphism. The reader will find all this (and much more than he needs) in
``modern'' algebra textbooks, for example the excellent books by
R. Godement (Houghton--Mifflin) and S. Lang (Addison--Wesley). This book
uses throughout the language and results of ``modern'' algebra. For the
convenience of the reader I have included discussions of the following
notions: integrality, algebraic extensions of fields, Galois theory,
Noetherian rings and modules, and rings of fractions. I have tried to give
short and clear accounts of these topics but without unnecessary
sophistication.

This book exists as a consequence of a course in ``Higher Mathematics''
taught at the University of Paris in 1965 and repeated in 1966. Course notes
taken by Alfred Vidal-Madjar, a student at the École Normale Supérieure, to
whom I wish gratefully to acknowledge my appreciation, have served as a
first draft. Certain parts have come from courses taught at the École
Normale Supérieure de Jeunes Filles and at the University of
Clermont-sur-Tiretaine. Finally the influence and advice of several
mathematicians have been valuable to me. Among them I want particularly to
thank the master of my generation, N. Bourbaki, who has had the kindness to
show me his unpublished manuscripts, also my friends Emil Grosswald,
Georges Poitou, Jean-Pierre Serre, and John Tate.

\vspace{1.2\baselineskip}
\begin{flushright}
Bourg-la-Reine, November 1966
\end{flushright}
\vspace{2.5\baselineskip}
\begin{flushright}\small\itshape
A NICOLE\\
qui a su créer autour de moi\\
une atmosphère favorable à ce livre
\end{flushright}

%% --------------- notations, definitions, and prerequisites ------------
\clearpage
\setcounter{page}{11}%
\fmhead{Notations, definitions, and prerequisites}%
       {NOTATIONS, DEFINITIONS AND PREREQUISITES}{notations}

We employ the usual notations from set theory: $\in,\subset,\cup,\cap$.
The complement of a subset $B$ of a set $A$ is denoted $A-B$. $\phi$ denotes
the empty set. The cardinality (or power, or number of elements) of a set
$A$ is written $\operatorname{card}(A)$; if $A$ is a group, we speak of the
\emph{order} of $A$.

We assume that the reader is acquainted with the notions of group, ring,
field and vector space, as well as with the elementary theory of vector
spaces (also called ``linear algebra''). In this book, with the exception of
Chapter V, § 7, ``ring'' (respectively, ``field'') means commutative ring
(respectively, field) with unit element.

Given a finite group $G$ and a subgroup $H$ of $G$, we recall that
$\operatorname{card}(H)$ divides $\operatorname{card}(G)$. The quotient
$\operatorname{card}(G)/\operatorname{card}(H)$ is called the index of $H$
in $G$ and is denoted $(G:H)$.

Given two subsets $A$ and $B$ of an additive group $G$, we write $A+B$ for
the set of sums $a+b$, $a\in A$, $b\in B$.

Given a ring $A$, we write $A[X]$ or $A[Y]$ (capital letter) for the ring of
polynomials in one variable over $A$; we write $A[X_{1},\ldots,X_{n}]$ for
the polynomials in $n$ variables and $A[[X]]$ for the ring of formal power
series.

By convention, a subring $A$ of a ring $B$ contains the unit element of $B$.
Given a ring $B$, a subring $A$ of $B$, and an element $x\in B$, we write
$A[x]$ for the subring of $B$ generated by $A$ and $x$, i.e.\ for the
intersection of all subrings of $B$ which contain $A$ and $x$; it is the set
of all sums of the form $a_{0}+a_{1}x+\cdots+a_{n}x^{n}$ ($a_{i}\in A$). We
write $A[x_{1},\ldots,x_{n}]$ for the subring of $B$ generated by $A$ and a
finite set $(x_{1},\ldots,x_{n})$ of elements of $B$.

A ring $A$ is called an \emph{integral domain} if $A$ contains more than one
element and if the product of any two non-zero elements of $A$ is not zero.

An ideal $\mathfrak{b}$ of a ring $A$ is a subgroup of the additive group of
$A$ such that $x\in\mathfrak{b}$ and $a\in A$ implies $ax\in\mathfrak{b}$.
The whole ring and the set consisting of the element $0$ alone (and denoted
$(0)$) are ideals, sometimes called ``trivial'' ideals. A field has no
non-trivial ideals and this fact distinguishes fields from other rings.
Given a set of elements $(b_{t})$ from a ring $A$, the intersection of all
ideals of $A$ containing the $b_{t}$'s is an ideal of $A$, called the ideal
generated by the $b_{t}$'s; it is the set of all elements of the form
$\sum a_{t}b_{t}$ with $a_{t}\in A$.

An ideal generated by a single element $b$ is called \emph{principal};
notation: $Ab$ or $(b)$.

Let $A$ be a ring and $\mathfrak{b}$ an ideal of $A$. The equivalence classes
$a+\mathfrak{b}$ ($a\in A$) form a ring, called the quotient ring of $A$ by
$\mathfrak{b}$ and denoted $A/\mathfrak{b}$. The ideals of $A/\mathfrak{b}$
are of the form $\mathfrak{b}'/\mathfrak{b}$ where $\mathfrak{b}'$ runs
through the set of ideals of $A$ which contain $\mathfrak{b}$. In order that
$A/\mathfrak{b}$ be a field it is necessary and sufficient that
$\mathfrak{b}$ be maximal among the ideals of $A$ distinct from $A$. We say
then that $\mathfrak{b}$ is a \emph{maximal ideal}. An ideal
$\mathfrak{p}$ is called \emph{prime} if $A/\mathfrak{p}$ is an integral
domain.

Let $A$ and $A'$ be rings with unit elements $e$ and $e'$. A homomorphism
$f:A\to A'$ is a mapping $f$ of $A$ to $A'$ such that:
\[
f(a+b)=f(a)+f(b),\quad f(ab)=f(a)f(b),\quad\text{and}\quad f(e)=e'.
\]

Let $A$ be a ring. An $A$-algebra is a pair consisting of a ring $B$ and a
homomorphism $\varphi:\mathbf{A}\to\mathbf{B}$. If $A$ is a field,
$\varphi$ is injective, and we often identify $A$ with its image
$\varphi(\mathbf{A})$ (which is a subring of $B$).

Given a field $L$ and a subfield $K$ of $L$, we call $L$ an
\emph{extension} of $K$.

Usually the unit element of a ring $A$ will be denoted by $1$.

The notion of module over a ring $A$ (or $A$-module) is the direct
generalization of the notion of vector space over a field. An $A$-module $M$
is an abelian group (the operation is addition) provided with a mapping
$A\times M\to M$ (written as multiplication) such that
$a(x+y)=ax+ay$, $(a+b)x=ax+bx$, $a(bx)=(ab)x$, $1x=x$
($a,b\in A$, $x,y\in M$). There are notions of submodule and of quotient
module. Given two $A$-modules, $M$ and $M'$, a homomorphism (or $A$-linear
mapping) of $M$ to $M'$ is a mapping $f:M\to M'$ such that
\[
f(x+y)=f(x)+f(y)\quad\text{and}\quad f(ax)=af(x)\quad(a\in A,\ x,y\in M).
\]

Given a homomorphism $f:\mathbf{X}\to\mathbf{X}'$ (of groups, rings, or
modules), we call the \emph{kernel} of $f$ and write
$\ker(f)$ for the inverse image under $f$ of the zero or identity element of
$\mathbf{X}'$. It is a normal subgroup (or an ideal, or a submodule) of
$\mathbf{X}$. In order that $f$ be injective it is necessary and sufficient
that $\ker(f)$ consist of only the zero or identity element of
$\mathbf{X}$. We call the \emph{image} of $f$ the subset $f(\mathbf{X})$
of $\mathbf{X}'$; it is a subgroup (or subring, or submodule) of
$\mathbf{X}'$.

Let $X$ and $X'$ be sets. A mapping $f$ of $X$ to $X'$ is frequently
denoted $f:X\to X'$. When a mapping $f:X\to X'$ is described by the value it
takes at an arbitrary element $x\in X$, we use the notation $x\mapsto f(x)$.
Thus the sine function, $\sin:\mathbf{R}\to\mathbf{R}$, can be defined by
\[
x\mapsto\sum_{n=0}^{\infty}(-1)^{n}\frac{x^{2n+1}}{(2n+1)!}.
\]

We shall employ the usual notations for the following mathematical objects:

\noindent\textbf{N}: set of natural numbers $(0,1,2,\ldots,n,\ldots)$\\
\hspace*{2.5em}(\textbf{N} for ``numbers'').\par
\vspace{2pt}
\noindent\textbf{Z}: ring of rational integers (natural numbers and their
negatives)\\
\hspace*{2.5em}(\textbf{Z} for ``Zahlen'').\par
\vspace{2pt}
\noindent\textbf{Q}: field of rational numbers (quotients of elements of
\textbf{Z})\\
\hspace*{2.5em}(\textbf{Q} for ``quotients'').\par
\vspace{2pt}
\noindent\textbf{R}: field of real numbers\\
\hspace*{2.5em}(\textbf{R} for ``reals'').\par
\vspace{2pt}
\noindent\textbf{C}: field of complex numbers\\
\hspace*{2.5em}(\textbf{C} for ``complexes'').\par
\vspace{2pt}
\noindent$\mathbf{F}_{q}$: finite field with $q$ elements\\
\hspace*{2.5em}(\textbf{F} for ``finite'' or ``field'')\par
